\begingroup\setbox0=\vbox{\section*{Tỉ lệ thuận và tỉ lệ nghịch}}\endgroup
\chapterhead{7}{Tỉ lệ thuận và tỉ lệ nghịch}

\shead{Hai đại lượng tỉ lệ}
\begin{knowbox}
\textbf{Tỉ lệ thuận.} Hai đại lượng tỉ lệ thuận với nhau nếu khi đại lượng này gấp lên (hoặc giảm đi) bao nhiêu lần thì đại lượng kia cũng gấp lên (hoặc giảm đi) bấy nhiêu lần. \emph{Ví dụ:} số mét vải và số tiền mua vải (khi giá không đổi).\par
\textbf{Tỉ lệ nghịch.} Hai đại lượng tỉ lệ nghịch với nhau nếu khi đại lượng này gấp lên bao nhiêu lần thì đại lượng kia giảm đi bấy nhiêu lần; \emph{tích} của chúng không đổi. \emph{Ví dụ:} số công nhân và số ngày cần để làm xong cùng một công việc.\par
\smallskip
Hai cách giải thường dùng:
\begin{itemize}[leftmargin=1.4em,itemsep=3pt]
\item \textbf{Phương pháp rút về đơn vị:} tìm giá trị ứng với $1$ đơn vị (1 mét, 1 công nhân, 1 ngày,\ldots), rồi nhân hoặc chia.
\item \textbf{Phương pháp tìm tỉ số:} tìm xem đại lượng thứ nhất gấp lên (hoặc giảm đi) bao nhiêu lần, rồi thay đổi đại lượng thứ hai tương ứng.
\end{itemize}
Với các bài toán ``công nhân -- ngày'', ta coi công việc là \emph{(số công nhân) $\times$ (số ngày)} ``ngày công''; số ngày công cần thiết không thay đổi.
\end{knowbox}
\gap
\begin{exbox}{\ExWord\ 1}
Mua $5$ m vải hết $450\,000$ đồng. Hỏi mua $15$ m vải cùng loại hết bao nhiêu tiền?
\guide
\textbf{Cách 1 (rút về đơn vị).} Mua $1$ m vải hết $450\,000:5=90\,000$ đồng, vậy mua $15$ m vải hết $90\,000\times15=1\,350\,000$ đồng.\par
\textbf{Cách 2 (tìm tỉ số).} $15$ m gấp $15:5=3$ lần $5$ m, nên số tiền cũng gấp $3$ lần: $450\,000\times3=1\,350\,000$ đồng.
\dapso{$1\,350\,000$ đồng.}
\end{exbox}
\gap
\begin{exbox}{\ExWord\ 2}
Một đội $12$ công nhân làm xong một đoạn đường trong $10$ ngày. Muốn làm xong đoạn đường đó trong $8$ ngày thì cần bao nhiêu công nhân? (Năng suất làm việc của mỗi người như nhau.)
\guide
Số công nhân và số ngày là hai đại lượng tỉ lệ nghịch.\par Một công nhân làm một mình sẽ cần $10\times12=120$ ngày, vậy công việc cần $120$ ngày công.\par
Để làm xong trong $8$ ngày cần $120:8=15$ công nhân.
\dapso{$15$ công nhân.}
\end{exbox}
\gap
\begin{exbox}{\ExWord\ 3}
Một bếp ăn dự trữ gạo đủ cho $120$ người ăn trong $20$ ngày. Sau $5$ ngày, có thêm $30$ người đến. Hỏi số gạo còn lại đủ ăn thêm bao nhiêu ngày nữa? (Mức ăn của mỗi người như nhau.)
\guide
Sau $5$ ngày, số gạo còn lại đủ cho $120$ người ăn trong $20-5=15$ ngày.\par
Số gạo đó đủ cho $1$ người ăn trong $15\times120=1\,800$ ngày.\par
Bây giờ có $120+30=150$ người, nên số gạo đủ ăn trong $1\,800:150=12$ ngày.
\dapso{$12$ ngày.}
\end{exbox}
\gap
\begin{exbox}{\ExWord\ 4}
$10$ công nhân sửa được $400$ m đường trong $5$ ngày. Hỏi $12$ công nhân sửa được bao nhiêu mét đường trong $6$ ngày? (Năng suất làm việc của mỗi người như nhau.)
\guide
Có ba đại lượng; ta rút về ``$1$ công nhân trong $1$ ngày''.\par
Trong $5$ ngày, $1$ công nhân sửa được $400:10=40$ m; trong $1$ ngày sửa được $40:5=8$ m.\par
$12$ công nhân trong $6$ ngày sửa được $8\times12\times6=576$ m.
\dapso{$576$ m.}
\end{exbox}
\dividerbar{Luyện tập}
\begin{probbox}{1}{\relax}
$6$ hộp bút như nhau đựng tất cả $72$ cái bút. Hỏi $15$ hộp như thế đựng bao nhiêu cái bút?
\end{probbox}
\probsol{
Mỗi hộp đựng $72:6=12$ cái bút, vậy $15$ hộp đựng $12\times15=180$ cái bút.
\dapso{$180$ cái bút.}}
\gap
\begin{probbox}{2}{\relax}
Một ô tô tiêu thụ hết $9$ lít xăng khi đi $120$ km. Hỏi ô tô đó đi $200$ km thì tiêu thụ hết bao nhiêu lít xăng?
\end{probbox}
\probsol{
Đi $40$ km hết $9:3=3$ lít (vì $120=40\times3$).\par $200$ km gấp $200:40=5$ lần $40$ km, nên tiêu thụ hết $3\times5=15$ lít.
\dapso{$15$ lít.}}
\gap
\begin{probbox}{3}{\relax}
May $15$ bộ quần áo hết $45$ m vải. Hỏi có $72$ m vải thì may được bao nhiêu bộ quần áo như thế?
\end{probbox}
\probsol{
Mỗi bộ cần $45:15=3$ m vải, vậy $72$ m vải may được $72:3=24$ bộ.
\dapso{$24$ bộ quần áo.}}
\gap
\begin{probbox}{4}{\relax}
Có một số tiền đủ mua $24$ quyển vở, giá $6\,000$ đồng một quyển. Hỏi với số tiền đó thì mua được bao nhiêu quyển vở loại giá $8\,000$ đồng một quyển?
\end{probbox}
\probsol{
Số tiền đó là $6\,000\times24=144\,000$ đồng, mua được $144\,000:8\,000=18$ quyển vở.
\dapso{$18$ quyển vở.}}
\gap
\begin{probbox}{5}{\relax}
$15$ công nhân làm xong một công việc trong $12$ ngày. Hỏi $20$ công nhân làm xong công việc đó trong bao nhiêu ngày? (Năng suất làm việc của mỗi người như nhau.)
\end{probbox}
\probsol{
Công việc cần $15\times12=180$ ngày công, vậy $20$ công nhân cần $180:20=9$ ngày.
\dapso{$9$ ngày.}}
\gap
\begin{probbox}{6}{\relax}
Một bếp ăn dự trữ gạo đủ cho $80$ người ăn trong $15$ ngày. Hỏi số gạo đó đủ cho $100$ người ăn trong bao nhiêu ngày?
\end{probbox}
\probsol{
Số gạo đó đủ cho $1$ người ăn trong $15\times80=1\,200$ ngày, vậy đủ cho $100$ người ăn trong $1\,200:100=12$ ngày.
\dapso{$12$ ngày.}}
\gap
\begin{probbox}{7}{\relax}
Một trại hè có $60$ học sinh, dự trữ thức ăn đủ dùng trong $30$ ngày. Sau $10$ ngày, có $20$ học sinh về sớm. Hỏi số thức ăn còn lại đủ cho số học sinh ở lại dùng thêm bao nhiêu ngày nữa?
\end{probbox}
\probsol{
Số thức ăn còn lại đủ cho $60$ học sinh dùng trong $20$ ngày, tức là\par $60\times20=1\,200$ suất ăn (mỗi suất cho một học sinh trong một ngày).\par Bây giờ còn $40$ học sinh, nên đủ dùng trong $1\,200:40=30$ ngày.
\dapso{$30$ ngày.}}
\gap
\begin{probbox}{8}{\relax}
$8$ người trồng được $96$ cây trong $3$ giờ. Hỏi $10$ người trồng được bao nhiêu cây trong $4$ giờ? (Năng suất làm việc của mỗi người như nhau.)
\end{probbox}
\probsol{
$1$ người trồng được $96:8:3=4$ cây mỗi giờ.\par $10$ người trong $4$ giờ: $4\times10\times4=160$ cây.
\dapso{$160$ cây.}}
\gap
\begin{probbox}{9}{\relax}
$3$ công nhân làm xong một công việc trong $8$ ngày. Hỏi cần thêm bao nhiêu công nhân để làm xong công việc đó sớm hơn $2$ ngày? (Năng suất làm việc của mỗi người như nhau.)
\end{probbox}
\probsol{
Công việc cần $3\times8=24$ ngày công.\par Trong $8-2=6$ ngày cần $24:6=4$ công nhân, tức là\par cần thêm $4-3=1$ người.
\dapso{$1$ công nhân.}}
\gap
\begin{probbox}{10}{\relax}
Cứ $100$ kg thóc thì xay được $62$ kg gạo. Hỏi $2\,500$ kg thóc thì xay được bao nhiêu ki-lô-gam gạo?
\end{probbox}
\probsol{
$2\,500$ kg gấp $25$ lần $100$ kg, nên số gạo là $62\times25=1\,550$ kg.
\dapso{$1\,550$ kg.}}
\gap
\begin{probbox}{11}{\relax}
$5$ máy bơm như nhau bơm được $360$ m$^3$ nước trong $4$ giờ. Hỏi $8$ máy bơm như thế bơm được bao nhiêu nước trong $3$ giờ?
\end{probbox}
\probsol{
$1$ máy trong $1$ giờ: $360:5:4=18$ m$^3$.\par $8$ máy trong $3$ giờ: $18\times8\times3=432$ m$^3$.
\dapso{$432$ m$^3$.}}
\gap
\begin{probbox}{12}{\relax}
Một xưởng có $45$ công nhân dự định làm xong một đơn hàng trong $20$ ngày. Sau $4$ ngày, có thêm $15$ công nhân đến làm cùng. Hỏi đơn hàng được hoàn thành sớm hơn dự định bao nhiêu ngày?
\end{probbox}
\probsol{
Phần việc còn lại là của $45$ công nhân làm trong $16$ ngày, tức là\par $45\times16=720$ ngày công.\par Với $60$ công nhân thì cần $720:60=12$ ngày.\par
Tổng cộng $4+12=16$ ngày, tức là sớm hơn $20-16=4$ ngày.
\dapso{$4$ ngày.}}
\gap
\begin{probbox}{13}{\relax}
Một bếp ăn dự trữ gạo đủ cho $60$ người ăn trong $24$ ngày. Sau $4$ ngày có thêm một số người đến, nên số gạo còn lại chỉ đủ ăn thêm $16$ ngày nữa. Hỏi có bao nhiêu người đến thêm?
\end{probbox}
\probsol{
Số gạo còn lại là $60\times20=1\,200$ suất ăn.\par Trong $16$ ngày, số gạo đó đủ cho $1\,200:16=75$ người.\par
Vậy có $75-60=15$ người đến thêm.
\dapso{$15$ người.}}
\gap
\begin{probbox}{14}{\relax}
Mua $4$ kg táo hết $120\,000$ đồng. Hỏi mua $7$ kg táo hết bao nhiêu tiền?
\end{probbox}
\probsol{
$1$ kg: $120\,000:4=30\,000$ đồng.\par $7$ kg: $30\,000\times7=210\,000$ đồng.
\dapso{$210\,000$ đồng.}}
\gap
\begin{probbox}{15}{\relax}
Mua $12$ quyển vở hết $96\,000$ đồng. Hỏi mua $20$ quyển hết bao nhiêu tiền?
\end{probbox}
\probsol{
Một quyển: $96\,000:12=8\,000$ đồng.\par $20$ quyển: $8\,000\times20=160\,000$ đồng.
\dapso{$160\,000$ đồng.}}
\gap
\begin{probbox}{16}{\relax}
Một ô tô đi được $135$ km trong $3$ giờ. Hỏi trong $5$ giờ ô tô đi được bao nhiêu ki-lô-mét? (Mỗi giờ đi được quãng đường như nhau.)
\end{probbox}
\probsol{
Một giờ: $135:3=45$ km.\par $5$ giờ: $45\times5=225$ km.
\dapso{$225$ km.}}
\gap
\begin{probbox}{17}{\relax}
$8$ công nhân làm xong một công việc trong $6$ ngày. Hỏi $12$ công nhân làm xong công việc đó trong bao nhiêu ngày? (Năng suất như nhau.)
\end{probbox}
\probsol{
Công việc cần $8\times6=48$ ngày công.\par $12$ công nhân cần $48:12=4$ ngày.
\dapso{$4$ ngày.}}
\gap
\begin{probbox}{18}{\relax}
Mỗi ngày đọc $15$ trang thì Ha đọc xong một quyển sách trong $12$ ngày. Hỏi nếu mỗi ngày đọc $20$ trang thì Ha đọc xong quyển sách đó trong bao nhiêu ngày?
\end{probbox}
\probsol{
Quyển sách có $15\times12=180$ trang.\par Số ngày: $180:20=9$.
\dapso{$9$ ngày.}}
\gap
\begin{probbox}{19}{\relax}
Một bếp ăn dự trữ gạo đủ cho $50$ người ăn trong $18$ ngày. Hỏi số gạo đó đủ cho $60$ người ăn trong bao nhiêu ngày?
\end{probbox}
\probsol{
Số gạo đủ cho $1$ người ăn trong $18\times50=900$ ngày.\par $60$ người: $900:60=15$ ngày.
\dapso{$15$ ngày.}}
\gap
\begin{probbox}{20}{\relax}
$6$ máy khâu như nhau may được $600$ cái áo trong $5$ ngày. Hỏi $10$ máy may được bao nhiêu cái áo trong $3$ ngày?
\end{probbox}
\probsol{
$1$ máy trong $1$ ngày: $600:6:5=20$ cái áo.\par $10$ máy trong $3$ ngày: $20\times10\times3=600$ cái áo.
\dapso{$600$ cái áo.}}
\gap
\begin{probbox}{21}{\relax}
Mua $6$ m dây điện hết $54\,000$ đồng. Hỏi với $360\,000$ đồng thì mua được bao nhiêu mét dây điện?
\end{probbox}
\probsol{
$1$ m giá $54\,000:6=9\,000$ đồng.\par $360\,000:9\,000=40$ m.
\dapso{$40$ m.}}
\gap
\begin{probbox}{22}{\relax}
Một vòi nước chảy được $60$ lít trong $5$ phút. Hỏi trong $1$ giờ vòi đó chảy được bao nhiêu lít nước?
\end{probbox}
\probsol{
$1$ giờ $=60$ phút, gấp $60:5=12$ lần $5$ phút.\par Lượng nước: $60\times12=720$ lít.
\dapso{$720$ lít.}}
\gap
\begin{probbox}{23}{\relax}
$6$ xe tải như nhau chở được $72$ tấn hàng. Hỏi $15$ xe tải như thế chở được bao nhiêu tấn hàng?
\end{probbox}
\probsol{
Mỗi xe: $72:6=12$ tấn.\par $15$ xe: $12\times15=180$ tấn.
\dapso{$180$ tấn.}}
\gap
\begin{probbox}{24}{\relax}
Một đội $20$ người dự định làm xong một công việc trong $15$ ngày. Hỏi cần thêm bao nhiêu người để làm xong trong $12$ ngày? (Năng suất như nhau.)
\end{probbox}
\probsol{
Công việc cần $20\times15=300$ ngày công.\par Trong $12$ ngày cần $300:12=25$ người.\par Vậy cần thêm $25-20=5$ người.
\dapso{$5$ người.}}
\gap
\begin{probbox}{25}{\relax}
$9$ người dự định làm xong một công việc trong $12$ ngày. Sau $4$ ngày, có $3$ người nghỉ việc. Hỏi công việc hoàn thành chậm hơn dự định bao nhiêu ngày?
\end{probbox}
\probsol{
Phần việc còn lại: $9\times(12-4)=72$ ngày công.\par $6$ người cần $72:6=12$ ngày.\par Tổng cộng: $4+12=16$ ngày, tức là chậm hơn $16-12=4$ ngày.
\dapso{$4$ ngày.}}
\gap
\begin{probbox}{26}{\relax}
Cứ $30$ kg thóc thì xay được $18$ kg gạo. Hỏi $1$ tấn thóc thì xay được bao nhiêu gạo?
\end{probbox}
\probsol{
$1$ tấn $=1\,000$ kg.\par $1$ kg thóc xay được $18:30=0{,}6$ kg gạo.\par $1\,000$ kg thóc xay được $0{,}6\times1\,000=600$ kg gạo.
\dapso{$600$ kg.}}
\gap
\begin{probbox}{27}{\relax}
$5$ bao gạo như nhau cân nặng $250$ kg. Hỏi $12$ bao gạo như thế cân nặng bao nhiêu?
\end{probbox}
\probsol{
Một bao: $250:5=50$ kg.\par $12$ bao: $50\times12=600$ kg.
\dapso{$600$ kg.}}
\gap
\begin{probbox}{28}{\relax}
Một trường trả $9\,000\,000$ đồng để mua $15$ bộ bàn ghế. Hỏi muốn mua $25$ bộ thì cần thêm bao nhiêu tiền nữa?
\end{probbox}
\probsol{
Một bộ: $9\,000\,000:15=600\,000$ đồng.\par $25$ bộ: $600\,000\times25=15\,000\,000$ đồng.\par Cần thêm: $15\,000\,000-9\,000\,000=6\,000\,000$ đồng.
\dapso{$6\,000\,000$ đồng.}}
\gap
\begin{probbox}{29}{\relax}
Mỗi chuyến chở $4$ tấn thì một xe tải cần chở $15$ chuyến. Hỏi nếu mỗi chuyến chở $6$ tấn thì cần bao nhiêu chuyến?
\end{probbox}
\probsol{
Số hàng: $4\times15=60$ tấn.\par Số chuyến: $60:6=10$.
\dapso{$10$ chuyến.}}
\gap
\begin{probbox}{30}{\relax}
$14$ người ăn hết $21$ kg gạo trong $3$ ngày. Hỏi $20$ người ăn hết bao nhiêu ki-lô-gam gạo trong $6$ ngày? (Mức ăn của mỗi người như nhau.)
\end{probbox}
\probsol{
$1$ người trong $1$ ngày: $21:14:3=0{,}5$ kg.\par $20$ người trong $6$ ngày: $0{,}5\times20\times6=60$ kg.
\dapso{$60$ kg.}}
\dividerbar{Thử thách -- Mở rộng}
\begin{probbox}{31}{\stars{1}}
Một đơn vị dự trữ gạo đủ ăn trong $30$ ngày. Sau $10$ ngày, có thêm $40$ người đến, nên số gạo còn lại chỉ đủ ăn thêm $15$ ngày nữa. Hỏi lúc đầu đơn vị có bao nhiêu người?
\end{probbox}
\probsol{
Số gạo còn lại đủ cho số người lúc đầu ăn trong $20$ ngày, hoặc cho số người lúc đầu cộng thêm $40$ người ăn trong $15$ ngày.\par
Vì $20:15=\dfrac43$ nên số người lúc sau bằng $\dfrac43$ số người lúc đầu; vậy $40$ người bằng $\dfrac13$ số người lúc đầu.\par
Số người lúc đầu: $40\times3=120$ người.
\dapso{$120$ người.}}
\gap
\begin{probbox}{32}{\stars{1}}
Một đội dự định làm xong một công trình trong $15$ ngày. Sau $5$ ngày, có $6$ công nhân chuyển đi, nên đội phải làm thêm $2$ ngày so với dự định. Hỏi lúc đầu đội có bao nhiêu công nhân?
\end{probbox}
\probsol{
Phần việc còn lại dự định làm trong $10$ ngày nhưng thực tế làm trong $12$ ngày.\par Số ngày gấp $\dfrac{12}{10}=\dfrac65$ lần, nên số công nhân còn lại bằng $\dfrac56$ số công nhân lúc đầu.\par
Vậy $6$ công nhân bằng $1-\dfrac56=\dfrac16$ số công nhân của đội, nên lúc đầu đội có $6\times6=36$ công nhân.
\dapso{$36$ công nhân.}}
\gap
\begin{probbox}{33}{\stars{1}}
$12$ người, mỗi ngày làm $8$ giờ, thì làm xong một công việc trong $10$ ngày. Hỏi $16$ người, mỗi ngày làm $6$ giờ, thì cần bao nhiêu ngày? (Năng suất làm việc của mỗi người như nhau.)
\end{probbox}
\probsol{
Công việc cần $12\times8\times10=960$ giờ công.\par Mỗi ngày $16$ người làm được $16\times6=96$ giờ công.\par
Số ngày cần: $960:96=10$.
\dapso{$10$ ngày.}}
\gap
\begin{probbox}{34}{\stars{1}}
Hai bánh xe răng cưa ăn khớp với nhau. Bánh lớn có $48$ răng, bánh nhỏ có $18$ răng. Khi bánh lớn quay được $15$ vòng thì bánh nhỏ quay được bao nhiêu vòng?
\end{probbox}
\probsol{
Hai bánh xe ăn khớp nên số răng đi qua chỗ tiếp xúc ở hai bánh là như nhau: $48\times15=720$.\par
Bánh nhỏ quay được $720:18=40$ vòng.
\note{Số răng và số vòng quay là hai đại lượng tỉ lệ nghịch.}
\dapso{$40$ vòng.}}
\gap
\begin{probbox}{35}{\stars{2}}
``Một con gà rưỡi đẻ một quả trứng rưỡi trong một ngày rưỡi.'' Hỏi $6$ con gà như thế đẻ được bao nhiêu quả trứng trong $6$ ngày?
\end{probbox}
\probsol{
$1{,}5$ con gà đẻ $1{,}5$ quả trứng trong $1{,}5$ ngày, nên $1{,}5$ con gà đẻ $1$ quả trứng trong $1$ ngày.\par
$6$ con gà gấp $4$ lần $1{,}5$ con gà, nên mỗi ngày đẻ $4$ quả trứng, và trong $6$ ngày đẻ $4\times6=24$ quả trứng.
\dapso{$24$ quả trứng.}}
\gap
\begin{probbox}{36}{\stars{2}}
Với một số tiền, cô Lan có thể mua $20$ kg gạo loại I hoặc $30$ kg gạo loại II. Nếu cô dùng số tiền đó mua cả hai loại với số ki-lô-gam bằng nhau thì mỗi loại mua được bao nhiêu ki-lô-gam?
\end{probbox}
\probsol{
$1$ kg gạo loại I giá bằng $\dfrac1{20}$ số tiền, $1$ kg gạo loại II giá bằng $\dfrac1{30}$ số tiền.\par Mỗi loại một ki-lô-gam giá bằng $\dfrac1{20}+\dfrac1{30}=\dfrac1{12}$ số tiền.\par
Vậy cô mua được mỗi loại $12$ kg.
\dapso{Mỗi loại: $12$ kg.}}
\gap
\begin{probbox}{37}{\stars{1}}
Một đội dự định làm xong một công trình trong $20$ ngày. Nếu có thêm $10$ công nhân thì đội làm xong sớm hơn $4$ ngày. Hỏi đội có bao nhiêu công nhân? (Năng suất như nhau.)
\end{probbox}
\probsol{
Số ngày giảm từ $20$ xuống $16$, tức là\par bằng $\dfrac{16}{20}=\dfrac45$ thời gian, nên số công nhân phải gấp lên $\dfrac54$ lần.\par Vậy $10$ công nhân bằng $\dfrac14$ số công nhân của đội.\par Số công nhân của đội: $10\times4=40$.
\dapso{$40$ công nhân.}}
\gap
\begin{probbox}{38}{\stars{1}}
Một bếp ăn dự trữ gạo đủ cho $100$ người ăn trong $30$ ngày. Sau $6$ ngày, có $20$ người chuyển đi. Hỏi số gạo còn lại đủ ăn trong bao nhiêu ngày?
\end{probbox}
\probsol{
Số gạo còn lại là $100\times24=2\,400$ suất ăn.\par Còn lại $80$ người: $2\,400:80=30$ ngày.
\dapso{$30$ ngày.}}
\gap
\begin{probbox}{39}{\stars{1}}
Một bếp ăn dự trữ gạo đủ ăn trong $24$ ngày. Sau $4$ ngày, có thêm $20$ người đến nên số gạo còn lại chỉ đủ ăn trong $15$ ngày. Hỏi lúc đầu có bao nhiêu người?
\end{probbox}
\probsol{
Số gạo còn lại đủ cho số người lúc đầu ăn trong $20$ ngày, hoặc cho số người đó cộng thêm $20$ người ăn trong $15$ ngày.\par $20:15=\dfrac43$, nên số người lúc sau bằng $\dfrac43$ số người lúc đầu; $20$ người mới đến bằng $\dfrac13$ số người lúc đầu.\par Lúc đầu: $60$ người.
\dapso{$60$ người.}}
\gap
\begin{probbox}{40}{\stars{1}}
$6$ vòi nước như nhau cùng chảy thì đầy bể trong $4$ giờ. Cả $6$ vòi cùng chảy trong $1$ giờ, sau đó khóa $2$ vòi. Hỏi các vòi còn lại phải chảy thêm bao lâu nữa thì đầy bể?
\end{probbox}
\probsol{
Để đầy bể cần $6\times4=24$ ``giờ vòi'' (lượng nước một vòi chảy trong một giờ).\par Giờ đầu chảy được $6$; còn lại $18$.\par $4$ vòi cần $18:4=4{,}5$ giờ $=4$ giờ $30$ phút.
\dapso{$4$ giờ $30$ phút.}}
\gap
\begin{probbox}{41}{\stars{1}}
Ba bánh xe răng cưa A, B, C ăn khớp nối tiếp nhau (A với B, B với C). Bánh A có $30$ răng, bánh B có $20$ răng và bánh C có $45$ răng. Khi bánh A quay được $12$ vòng thì bánh C quay được bao nhiêu vòng?
\end{probbox}
\probsol{
Số răng đi qua mỗi chỗ tiếp xúc là như nhau: $30\times12=360$.\par Bánh B quay $360:20=18$ vòng; bánh C cũng có $360$ răng đi qua nên quay $360:45=8$ vòng.
\dapso{$8$ vòng.}}
\gap
\begin{probbox}{42}{\stars{1}}
Một chiếc đồng hồ mỗi ngày chạy nhanh $3$ phút. Hỏi sau bao nhiêu ngày thì đồng hồ chạy nhanh $1$ giờ?
\end{probbox}
\probsol{
$1$ giờ $=60$ phút.\par Số ngày: $60:3=20$.
\dapso{$20$ ngày.}}
\gap
\begin{probbox}{43}{\stars{1}}
Đi bộ với vận tốc $4$ km/giờ thì Mai đến trường hết $30$ phút. Hỏi nếu đi xe đạp với vận tốc $12$ km/giờ thì Mai đến trường hết bao lâu?
\end{probbox}
\probsol{
Quãng đường không đổi; vận tốc gấp $12:4=3$ lần thì thời gian giảm đi bấy nhiêu lần, tức là còn một phần ba.\par Thời gian: $30:3=10$ phút.
\dapso{$10$ phút.}}
\gap
\begin{probbox}{44}{\stars{1}}
$4$ người, mỗi ngày làm $8$ giờ, thì làm xong một công việc trong $6$ ngày. Hỏi $6$ người, mỗi ngày làm $4$ giờ, thì làm xong công việc đó trong bao nhiêu ngày?
\end{probbox}
\probsol{
Công việc cần $4\times8\times6=192$ giờ công.\par Mỗi ngày $6$ người làm được $6\times4=24$ giờ công.\par Số ngày: $192:24=8$.
\dapso{$8$ ngày.}}
\gap
\begin{probbox}{45}{\stars{1}}
Có $18$ xe tải, mỗi xe chở $45$ bao hàng. Nếu chỉ dùng $15$ xe để chở hết số hàng đó thì mỗi xe phải chở thêm bao nhiêu bao?
\end{probbox}
\probsol{
Số bao: $45\times18=810$.\par Mỗi xe: $810:15=54$; tức là nhiều hơn $54-45=9$ bao.
\dapso{$9$ bao.}}
\gap
\begin{probbox}{46}{\stars{2}}
``Ba con mèo bắt được ba con chuột trong ba phút.'' Hỏi $100$ con mèo bắt $100$ con chuột hết bao lâu? (Các con mèo bắt nhanh như nhau.)
\end{probbox}
\probsol{
Ba con mèo bắt được ba con chuột trong $3$ phút, nên mỗi con mèo bắt $1$ con chuột trong $3$ phút.\par $100$ con mèo, mỗi con bắt $1$ con chuột, cũng hết $3$ phút.
\dapso{$3$ phút.}}
\gap
\begin{probbox}{47}{\stars{2}}
Một số tiền mua được $15$ kg táo hoặc $10$ kg lê. Sau khi mua $6$ kg táo, số tiền còn lại mua được bao nhiêu ki-lô-gam lê?
\end{probbox}
\probsol{
$1$ kg táo giá bằng $\dfrac1{15}$ số tiền; $6$ kg táo giá bằng $\dfrac6{15}=\dfrac25$ số tiền.\par Số tiền còn lại $\dfrac35$ mua được $10\times\dfrac35=6$ kg lê.
\dapso{$6$ kg lê.}}
\gap
\begin{probbox}{48}{\stars{2}}
Mua $3$ kg cam và $2$ kg táo hết $175\,000$ đồng, biết giá $1$ kg táo gấp đôi giá $1$ kg cam. Tìm giá tiền mỗi loại.
\end{probbox}
\probsol{
$2$ kg táo có giá bằng $4$ kg cam, nên số tiền phải trả bằng giá của $3+4=7$ kg cam.\par Cam: $175\,000:7=25\,000$ đồng/kg; táo: $50\,000$ đồng/kg.
\dapso{Cam $25\,000$, táo $50\,000$ đồng/kg.}}
\gap
\begin{probbox}{49}{\stars{2}}
Một đội $12$ người dự định đắp $360$ m đê trong $6$ ngày. Sau $2$ ngày, có thêm $4$ người đến làm cùng. Hỏi đội hoàn thành sớm hơn dự định bao nhiêu ngày?
\end{probbox}
\probsol{
Mỗi người mỗi ngày đắp được $360:12:6=5$ m.\par Trong $2$ ngày đội đắp được $12\times5\times2=120$ m; còn lại $240$ m.\par $16$ người mỗi ngày đắp được $80$ m, cần $240:80=3$ ngày.\par Tổng cộng $5$ ngày: sớm hơn $1$ ngày.
\dapso{$1$ ngày.}}
\gap
\begin{probbox}{50}{\stars{2}}
Cùng một số tiền mua được $8$ cái bút hoặc $12$ quyển vở. Mỗi cái bút đắt hơn mỗi quyển vở $2\,000$ đồng. Hỏi số tiền đó là bao nhiêu?
\end{probbox}
\probsol{
Mua $8$ cái bút thay cho $8$ quyển vở thì tốn thêm $8\times2\,000=16\,000$ đồng.\par Số tiền thêm đó bằng giá $12-8=4$ quyển vở, nên mỗi quyển vở giá $4\,000$ đồng.\par Số tiền: $4\,000\times12=48\,000$ đồng.
\dapso{$48\,000$ đồng.}}
\clearpage
\titlehead
\begingroup\setbox0=\vbox{\section*{Bài toán công việc}}\endgroup
\chapterhead{8}{Bài toán công việc}

\shead{Coi toàn bộ công việc là một đơn vị}
\begin{knowbox}
\begin{itemize}[leftmargin=1.4em,itemsep=3pt]
\item Coi \textbf{toàn bộ công việc} (hoặc \textbf{bể đầy nước}) là $1$ đơn vị.
\item Nếu một người làm xong công việc trong $n$ giờ thì \textbf{mỗi giờ} người đó làm được $\dfrac1n$ công việc.
\item Khi nhiều người \textbf{cùng làm}, phần việc làm được trong $1$ giờ bằng \emph{tổng} các phần việc của từng người.
\item \textbf{Vòi tháo nước ra} làm bể vơi đi, nên phần của nó được \emph{trừ đi}.
\item Thời gian cần $=$ (phần việc phải làm) $:$ (phần việc làm được trong $1$ giờ).
\end{itemize}
\end{knowbox}
\gap
\begin{exbox}{\ExWord\ 1}
Người thứ nhất làm một mình xong một công việc trong $6$ giờ, người thứ hai làm một mình xong trong $3$ giờ. Hỏi nếu hai người cùng làm thì sau bao lâu sẽ xong công việc?
\guide
Mỗi giờ người thứ nhất làm được $\dfrac16$ công việc, người thứ hai làm được $\dfrac13$ công việc.\par
Mỗi giờ cả hai người làm được $\dfrac16+\dfrac13=\dfrac12$ công việc.\par
Thời gian hai người cùng làm: $1:\dfrac12=2$ giờ.
\dapso{$2$ giờ.}
\end{exbox}
\gap
\begin{exbox}{\ExWord\ 2}
Hai người cùng làm thì xong một công việc trong $12$ ngày. Nếu người thứ nhất làm một mình thì xong trong $20$ ngày. Hỏi người thứ hai làm một mình thì xong trong bao lâu?
\guide
Mỗi ngày cả hai người làm được $\dfrac1{12}$ công việc, người thứ nhất làm được $\dfrac1{20}$ công việc.\par
Vậy mỗi ngày người thứ hai làm được $\dfrac1{12}-\dfrac1{20}=\dfrac5{60}-\dfrac3{60}=\dfrac1{30}$ công việc, nên làm một mình thì xong trong $30$ ngày.
\dapso{$30$ ngày.}
\end{exbox}
\gap
\begin{exbox}{\ExWord\ 3}
Vòi I chảy đầy bể trong $4$ giờ, vòi II chảy đầy bể trong $6$ giờ. Bể đang cạn; cả hai vòi cùng chảy trong $1$ giờ, sau đó khóa vòi I. Hỏi vòi II phải chảy thêm bao lâu nữa thì đầy bể?
\guide
Trong $1$ giờ cả hai vòi chảy được $\dfrac14+\dfrac16=\dfrac5{12}$ bể, còn lại $1-\dfrac5{12}=\dfrac7{12}$ bể.\par
Vòi II cần $\dfrac7{12}:\dfrac16=\dfrac72$ giờ $=3$ giờ $30$ phút.
\dapso{$3$ giờ $30$ phút.}
\end{exbox}
\gap
\begin{exbox}{\ExWord\ 4}
Một vòi chảy vào đầy bể trong $3$ giờ; một vòi tháo nước ra hết bể đầy trong $5$ giờ. Bể đang cạn, mở cả hai vòi cùng lúc thì sau bao lâu bể đầy?
\guide
Mỗi giờ lượng nước trong bể tăng thêm $\dfrac13-\dfrac15=\dfrac2{15}$ bể.\par
Thời gian để đầy bể: $1:\dfrac2{15}=\dfrac{15}2$ giờ $=7$ giờ $30$ phút.
\dapso{$7$ giờ $30$ phút.}
\end{exbox}
\dividerbar{Luyện tập}
\begin{probbox}{1}{\relax}
An quét xong sân trường trong $30$ phút, Binh quét xong trong $20$ phút. Hỏi hai bạn cùng quét thì xong trong bao lâu?
\end{probbox}
\probsol{
Mỗi phút cả hai bạn quét được $\dfrac1{30}+\dfrac1{20}=\dfrac5{60}=\dfrac1{12}$ sân, nên cùng quét thì xong trong $12$ phút.
\dapso{$12$ phút.}}
\gap
\begin{probbox}{2}{\relax}
Vòi I chảy đầy bể trong $2$ giờ, vòi II chảy đầy bể trong $3$ giờ. Bể đang cạn, mở cả hai vòi thì sau bao lâu bể đầy?
\end{probbox}
\probsol{
Mỗi giờ cả hai vòi chảy được $\dfrac12+\dfrac13=\dfrac56$ bể.\par Thời gian: $1:\dfrac56=\dfrac65$ giờ $=1$ giờ $12$ phút.
\dapso{$1$ giờ $12$ phút.}}
\gap
\begin{probbox}{3}{\relax}
Hai công nhân cùng làm thì xong một công việc trong $6$ giờ. Người thứ nhất làm một mình thì xong trong $10$ giờ. Hỏi người thứ hai làm một mình thì xong trong bao lâu?
\end{probbox}
\probsol{
Mỗi giờ người thứ hai làm được $\dfrac16-\dfrac1{10}=\dfrac1{15}$ công việc, nên làm một mình thì xong trong $15$ giờ.
\dapso{$15$ giờ.}}
\gap
\begin{probbox}{4}{\relax}
Hai người cùng gặt xong một thửa ruộng trong $2$ giờ. Người thứ nhất gặt một mình thì xong trong $3$ giờ. Hỏi mỗi giờ người thứ hai gặt được mấy phần thửa ruộng?
\end{probbox}
\probsol{
Mỗi giờ người thứ hai gặt được $\dfrac12-\dfrac13=\dfrac16$ thửa ruộng.
\dapso{$\dfrac16$ thửa ruộng.}}
\gap
\begin{probbox}{5}{\relax}
Ba bạn cùng dọn thư viện. Nếu làm một mình thì Mai cần $4$ giờ, Hung cần $6$ giờ và Lan cần $12$ giờ. Hỏi cả ba bạn cùng làm thì xong trong bao lâu?
\end{probbox}
\probsol{
Mỗi giờ cả ba bạn làm được $\dfrac14+\dfrac16+\dfrac1{12}=\dfrac{3+2+1}{12}=\dfrac12$ công việc, nên xong trong $2$ giờ.
\dapso{$2$ giờ.}}
\gap
\begin{probbox}{6}{\relax}
Một vòi chảy vào đầy bể trong $4$ giờ; một vòi tháo nước ra hết bể đầy trong $6$ giờ. Bể đang cạn, mở cả hai vòi thì sau bao lâu bể đầy?
\end{probbox}
\probsol{
Mỗi giờ lượng nước trong bể tăng thêm $\dfrac14-\dfrac16=\dfrac1{12}$ bể, nên bể đầy sau $12$ giờ.
\dapso{$12$ giờ.}}
\gap
\begin{probbox}{7}{\relax}
Một bể đã có nước bằng $\dfrac13$ bể. Một vòi chảy đầy bể cạn trong $6$ giờ. Hỏi vòi đó phải chảy bao lâu nữa thì đầy bể?
\end{probbox}
\probsol{
Phần bể còn phải chảy: $1-\dfrac13=\dfrac23$.\par Thời gian: $\dfrac23:\dfrac16=4$ giờ.
\dapso{$4$ giờ.}}
\gap
\begin{probbox}{8}{\relax}
Vòi I chảy đầy bể trong $5$ giờ, vòi II chảy đầy bể trong $10$ giờ. Bể đang cạn, vòi I chảy trong $2$ giờ thì khóa lại và mở vòi II. Hỏi vòi II phải chảy bao lâu nữa thì đầy bể?
\end{probbox}
\probsol{
Vòi I chảy được $\dfrac15\times2=\dfrac25$ bể, còn lại $\dfrac35$ bể.\par Vòi II cần $\dfrac35:\dfrac1{10}=6$ giờ.
\dapso{$6$ giờ.}}
\gap
\begin{probbox}{9}{\relax}
Đội A làm một mình xong một công việc trong $9$ ngày, đội B làm một mình xong trong $18$ ngày. Đội A làm một mình $3$ ngày, sau đó đội B đến làm cùng. Hỏi hai đội cùng làm thêm bao nhiêu ngày nữa thì xong?
\end{probbox}
\probsol{
Trong $3$ ngày đội A làm được $\dfrac39=\dfrac13$ công việc, còn lại $\dfrac23$ công việc.\par
Mỗi ngày hai đội cùng làm được $\dfrac19+\dfrac1{18}=\dfrac16$ công việc, nên cần $\dfrac23:\dfrac16=4$ ngày.
\dapso{$4$ ngày.}}
\gap
\begin{probbox}{10}{\relax}
Ông Ba làm một mình xong một công việc trong $10$ ngày, ông Tu làm một mình xong trong $15$ ngày. Hai ông cùng làm $4$ ngày thì ông Ba nghỉ. Hỏi ông Tu phải làm bao lâu nữa thì xong?
\end{probbox}
\probsol{
Mỗi ngày hai ông cùng làm được $\dfrac1{10}+\dfrac1{15}=\dfrac16$ công việc, nên trong $4$ ngày làm được $\dfrac46=\dfrac23$ công việc, còn lại $\dfrac13$.\par
Ông Tu cần $\dfrac13:\dfrac1{15}=5$ ngày.
\dapso{$5$ ngày.}}
\gap
\begin{probbox}{11}{\relax}
Hai người cùng làm thì xong một công việc trong $4$ giờ; người thứ nhất làm một mình thì xong trong $6$ giờ. Hai người cùng làm $2$ giờ thì người thứ nhất nghỉ. Hỏi người thứ hai phải làm thêm bao lâu nữa?
\end{probbox}
\probsol{
Mỗi giờ người thứ hai làm được $\dfrac14-\dfrac16=\dfrac1{12}$ công việc.\par Trong $2$ giờ cùng làm, hai người làm được $\dfrac24=\dfrac12$ công việc.\par
Người thứ hai cần làm thêm $\dfrac12:\dfrac1{12}=6$ giờ.
\dapso{$6$ giờ.}}
\gap
\begin{probbox}{12}{\relax}
Hai người cùng làm một công việc và được trả $1\,800\,000$ đồng. Nếu làm một mình thì người thứ nhất cần $6$ ngày, người thứ hai cần $3$ ngày. Nếu chia tiền theo phần việc mỗi người làm được thì mỗi người nhận được bao nhiêu tiền?
\end{probbox}
\probsol{
Mỗi ngày người thứ nhất làm được $\dfrac16$ công việc, người thứ hai làm được $\dfrac13=\dfrac26$ công việc, nên tỉ số phần việc của hai người là $1:2$.\par
Người thứ nhất: $1\,800\,000:3=600\,000$ đồng; người thứ hai: $1\,200\,000$ đồng.
\dapso{$600\,000$ đồng và $1\,200\,000$ đồng.}}
\gap
\begin{probbox}{13}{\relax}
Một người làm xong một công việc trong $4$ giờ, một người khác làm xong trong $12$ giờ. Hỏi hai người cùng làm thì xong trong bao lâu?
\end{probbox}
\probsol{
Mỗi giờ hai người cùng làm được: $\dfrac14+\dfrac1{12}=\dfrac4{12}=\dfrac13$ công việc.\par Hai người cùng làm thì xong trong $3$ giờ.
\dapso{$3$ giờ.}}
\gap
\begin{probbox}{14}{\relax}
Vòi I chảy đầy bể trong $5$ giờ, vòi II chảy đầy bể trong $20$ giờ. Bể đang cạn, mở cả hai vòi thì sau bao lâu bể đầy?
\end{probbox}
\probsol{
Mỗi giờ: $\dfrac15+\dfrac1{20}=\dfrac5{20}=\dfrac14$ bể.\par Bể đầy sau $4$ giờ.
\dapso{$4$ giờ.}}
\gap
\begin{probbox}{15}{\relax}
Hai người cùng làm thì xong một công việc trong $8$ ngày; người thứ nhất làm một mình thì xong trong $12$ ngày. Hỏi người thứ hai làm một mình thì xong trong bao lâu?
\end{probbox}
\probsol{
Mỗi ngày người thứ hai làm được $\dfrac18-\dfrac1{12}=\dfrac3{24}-\dfrac2{24}=\dfrac1{24}$ công việc.\par Làm một mình: $24$ ngày.
\dapso{$24$ ngày.}}
\gap
\begin{probbox}{16}{\relax}
Hai vòi cùng chảy thì đầy bể trong $2$ giờ; vòi I chảy một mình thì đầy bể trong $3$ giờ. Hỏi vòi II chảy một mình thì đầy bể trong bao lâu?
\end{probbox}
\probsol{
Mỗi giờ vòi II chảy được $\dfrac12-\dfrac13=\dfrac16$ bể.\par Chảy một mình: $6$ giờ.
\dapso{$6$ giờ.}}
\gap
\begin{probbox}{17}{\relax}
Ba công nhân, nếu làm một mình thì lần lượt cần $6$, $10$ và $15$ ngày. Hỏi cả ba người cùng làm thì xong trong bao lâu?
\end{probbox}
\probsol{
Mỗi ngày: $\dfrac16+\dfrac1{10}+\dfrac1{15}=\dfrac{5+3+2}{30}=\dfrac13$ công việc.\par Cả ba người: $3$ ngày.
\dapso{$3$ ngày.}}
\gap
\begin{probbox}{18}{\relax}
Một vòi chảy vào đầy bể trong $6$ giờ; một vòi tháo nước ra hết bể đầy trong $9$ giờ. Bể đang cạn, mở cả hai vòi thì sau bao lâu bể đầy?
\end{probbox}
\probsol{
Mỗi giờ: $\dfrac16-\dfrac19=\dfrac3{18}-\dfrac2{18}=\dfrac1{18}$ bể.\par Bể đầy sau $18$ giờ.
\dapso{$18$ giờ.}}
\gap
\begin{probbox}{19}{\relax}
Một bể đã có nước bằng $\dfrac14$ bể. Một vòi chảy đầy bể cạn trong $8$ giờ. Hỏi vòi đó cần chảy bao lâu nữa thì đầy bể?
\end{probbox}
\probsol{
Phần bể còn phải chảy: $\dfrac34$ bể.\par Thời gian: $\dfrac34:\dfrac18=6$ giờ.
\dapso{$6$ giờ.}}
\gap
\begin{probbox}{20}{\relax}
Một bể đang chứa nước bằng $\dfrac23$ bể. Một vòi tháo nước ra hết bể đầy trong $6$ giờ. Hỏi tháo hết nước trong bể này mất bao lâu?
\end{probbox}
\probsol{
Mỗi giờ vòi tháo được $\dfrac16$ bể.\par Thời gian: $\dfrac23:\dfrac16=4$ giờ.
\dapso{$4$ giờ.}}
\gap
\begin{probbox}{21}{\relax}
Đội A làm một mình thì cần $12$ ngày, đội B cần $6$ ngày. Đội A làm một mình $3$ ngày, sau đó đội B đến làm cùng. Hỏi hai đội phải làm thêm bao nhiêu ngày nữa?
\end{probbox}
\probsol{
Trong $3$ ngày đội A làm được $\dfrac3{12}=\dfrac14$ công việc; còn lại $\dfrac34$ công việc.\par Mỗi ngày hai đội cùng làm được $\dfrac1{12}+\dfrac16=\dfrac14$ công việc.\par Thời gian: $\dfrac34:\dfrac14=3$ ngày.
\dapso{$3$ ngày.}}
\gap
\begin{probbox}{22}{\relax}
Vòi I chảy đầy bể trong $4$ giờ, vòi II trong $12$ giờ. Bể đang cạn, vòi I chảy $1$ giờ thì mở thêm vòi II. Hỏi sau bao lâu nữa bể đầy?
\end{probbox}
\probsol{
Sau $1$ giờ vòi I chảy được $\dfrac14$ bể; còn lại $\dfrac34$ bể.\par Mỗi giờ cả hai vòi chảy được: $\dfrac14+\dfrac1{12}=\dfrac13$ bể.\par Thời gian: $\dfrac34:\dfrac13=\dfrac94$ giờ $=2$ giờ $15$ phút.
\dapso{$2$ giờ $15$ phút.}}
\gap
\begin{probbox}{23}{\relax}
Hai người cùng làm được $\dfrac12$ công việc trong $3$ ngày. Người thứ nhất làm một mình thì xong trong $10$ ngày. Hỏi người thứ hai làm một mình thì xong trong bao lâu?
\end{probbox}
\probsol{
Mỗi ngày hai người cùng làm được $\dfrac12:3=\dfrac16$ công việc.\par Mỗi ngày người thứ hai làm được $\dfrac16-\dfrac1{10}=\dfrac1{15}$ công việc.\par Làm một mình: $15$ ngày.
\dapso{$15$ ngày.}}
\gap
\begin{probbox}{24}{\relax}
An làm được $\dfrac13$ công việc trong $2$ giờ; Binh làm được $\dfrac12$ công việc đó trong $6$ giờ. Hỏi hai bạn cùng làm thì xong công việc trong bao lâu?
\end{probbox}
\probsol{
An làm cả công việc thì cần $6$ giờ; Binh cần $12$ giờ.\par Mỗi giờ cùng làm được: $\dfrac16+\dfrac1{12}=\dfrac14$ công việc.\par Vậy cần $4$ giờ.
\dapso{$4$ giờ.}}
\gap
\begin{probbox}{25}{\relax}
Mỗi giờ vòi I chảy được $\dfrac16$ bể, vòi II chảy được $\dfrac13$ bể. Hỏi bể đang cạn, cả hai vòi cùng chảy trong $1$ giờ $20$ phút thì được mấy phần bể?
\end{probbox}
\probsol{
$1$ giờ $20$ phút $=\dfrac43$ giờ.\par Mỗi giờ cả hai vòi chảy được: $\dfrac16+\dfrac13=\dfrac12$ bể.\par Trong $\dfrac43$ giờ: $\dfrac12\times\dfrac43=\dfrac23$.
\dapso{$\dfrac23$ bể.}}
\gap
\begin{probbox}{26}{\relax}
Hai đội cùng sửa một con đường thì xong trong $6$ ngày; đội I làm một mình thì xong trong $10$ ngày. Hai đội cùng làm $2$ ngày thì đội I chuyển đi. Hỏi đội II phải làm bao lâu nữa thì xong?
\end{probbox}
\probsol{
Mỗi ngày đội II làm được $\dfrac16-\dfrac1{10}=\dfrac1{15}$ công việc.\par Sau $2$ ngày cùng làm, hai đội làm được $\dfrac26=\dfrac13$ công việc; còn lại $\dfrac23$.\par Đội II cần $\dfrac23:\dfrac1{15}=10$ ngày.
\dapso{$10$ ngày.}}
\gap
\begin{probbox}{27}{\relax}
Một vòi chảy đầy bể trong $3$ giờ, nhưng vì bể bị rò nên phải mất $4$ giờ mới đầy. Hỏi nếu bể đầy thì chỗ rò làm chảy hết nước trong bao lâu?
\end{probbox}
\probsol{
Mỗi giờ chỗ rò làm mất $\dfrac13-\dfrac14=\dfrac1{12}$ bể.\par Chỗ rò làm chảy hết bể nước trong $12$ giờ.
\dapso{$12$ giờ.}}
\gap
\begin{probbox}{28}{\relax}
Hai công nhân cùng làm một công việc và được chia nhau $2\,000\,000$ đồng. Nếu làm một mình thì người thứ nhất cần $4$ ngày, người thứ hai cần $6$ ngày. Chia tiền theo phần việc làm được thì mỗi người nhận được bao nhiêu?
\end{probbox}
\probsol{
Mỗi ngày người thứ nhất làm được $\dfrac14=\dfrac3{12}$ công việc, người thứ hai làm được $\dfrac16=\dfrac2{12}$ công việc; tỉ số phần việc là $3:2$.\par Người thứ nhất: $2\,000\,000:5\times3=1\,200\,000$ đồng; người thứ hai: $800\,000$ đồng.
\dapso{$1\,200\,000$ đồng và $800\,000$ đồng.}}
\gap
\begin{probbox}{29}{\relax}
Vòi I chảy đầy bể trong $30$ phút, vòi II trong $45$ phút. Bể đang cạn, cả hai vòi cùng chảy thì sau bao nhiêu phút bể đầy?
\end{probbox}
\probsol{
Mỗi phút: $\dfrac1{30}+\dfrac1{45}=\dfrac3{90}+\dfrac2{90}=\dfrac1{18}$ bể.\par Bể đầy sau $18$ phút.
\dapso{$18$ phút.}}
\gap
\begin{probbox}{30}{\relax}
Hai máy cày cùng cày xong một cánh đồng trong $4$ ngày; máy I cày một mình thì xong trong $6$ ngày. Sau $2$ ngày cùng cày, máy I bị hỏng. Hỏi máy II phải cày bao lâu nữa thì xong?
\end{probbox}
\probsol{
Mỗi ngày máy II cày được $\dfrac14-\dfrac16=\dfrac1{12}$ cánh đồng.\par Sau $2$ ngày cùng cày, đã xong một nửa cánh đồng.\par Máy II cần $\dfrac12:\dfrac1{12}=6$ ngày.
\dapso{$6$ ngày.}}
\dividerbar{Thử thách -- Mở rộng}
\begin{probbox}{31}{\stars{1}}
Một vòi chảy đầy bể (bể không bị rò) trong $2$ giờ. Vì bể bị rò nên bây giờ vòi phải chảy $2$ giờ $30$ phút mới đầy bể. Hỏi nếu bể đầy thì chỗ rò làm chảy hết nước trong bao lâu?
\end{probbox}
\probsol{
$2$ giờ $30$ phút $=\dfrac52$ giờ.\par Mỗi giờ vòi chảy được $\dfrac12$ bể; khi bể bị rò, mỗi giờ nước trong bể chỉ tăng thêm $\dfrac25$ bể.\par
Mỗi giờ chỗ rò làm mất $\dfrac12-\dfrac25=\dfrac1{10}$ bể, nên làm chảy hết bể đầy trong $10$ giờ.
\dapso{$10$ giờ.}}
\gap
\begin{probbox}{32}{\stars{1}}
Một bể đang cạn. Vòi I chảy đầy bể trong $3$ giờ, vòi II trong $6$ giờ; vòi III tháo hết bể đầy trong $4$ giờ. Mở cả ba vòi thì sau bao lâu bể đầy?
\end{probbox}
\probsol{
Mỗi giờ nước trong bể tăng thêm $\dfrac13+\dfrac16-\dfrac14=\dfrac{4+2-3}{12}=\dfrac14$ bể, nên bể đầy sau $4$ giờ.
\dapso{$4$ giờ.}}
\gap
\begin{probbox}{33}{\stars{1}}
Có ba công nhân: nếu làm một mình thì A xong công việc trong $6$ ngày, B trong $12$ ngày và C trong $12$ ngày. A và B cùng làm $2$ ngày, sau đó C đến làm cùng. Hỏi cả ba người phải làm thêm bao lâu nữa?
\end{probbox}
\probsol{
Mỗi ngày A và B làm được $\dfrac16+\dfrac1{12}=\dfrac14$ công việc, nên trong $2$ ngày làm được một nửa công việc.\par
Mỗi ngày cả ba người làm được $\dfrac14+\dfrac1{12}=\dfrac13$ công việc, nên cần $\dfrac12:\dfrac13=\dfrac32$ ngày, tức là một ngày rưỡi.
\dapso{$1{,}5$ ngày.}}
\gap
\begin{probbox}{34}{\stars{1}}
Người thợ cả làm một mình xong một công việc trong $12$ giờ. Người thợ cả làm $3$ giờ, sau đó người thợ phụ làm $4$ giờ, thì cả hai làm được $\dfrac7{12}$ công việc. Hỏi nếu hai người cùng làm từ đầu thì xong công việc trong bao lâu?
\end{probbox}
\probsol{
Trong $3$ giờ người thợ cả làm được $\dfrac3{12}$ công việc.\par $4$ giờ của người thợ phụ làm được $\dfrac7{12}-\dfrac3{12}=\dfrac4{12}$ công việc, nên mỗi giờ người thợ phụ làm được $\dfrac1{12}$ công việc.\par
Mỗi giờ cùng làm được: $\dfrac1{12}+\dfrac1{12}=\dfrac16$ công việc, vậy cần $6$ giờ.
\dapso{$6$ giờ.}}
\gap
\begin{probbox}{35}{\stars{2}}
A và B cùng làm thì xong một công việc trong $12$ ngày; B và C cùng làm thì xong trong $15$ ngày; A và C cùng làm thì xong trong $20$ ngày.
\begin{cols2}
\pa{a} Cả ba người cùng làm thì xong trong bao lâu? & \pa{b} Mỗi người làm một mình thì xong trong bao lâu?
\end{cols2}
\end{probbox}
\probsol{
\pa{a} Cộng phần việc của ba cặp: $\dfrac1{12}+\dfrac1{15}+\dfrac1{20}=\dfrac{5+4+3}{60}=\dfrac15$.\par Mỗi người được tính hai lần, nên mỗi ngày cả ba người làm được $\dfrac15:2=\dfrac1{10}$ công việc và xong trong $10$ ngày.\par
\pa{b} A: $\dfrac1{10}-\dfrac1{15}=\dfrac1{30}$ ($30$ ngày); B: $\dfrac1{10}-\dfrac1{20}=\dfrac1{20}$ ($20$ ngày); C: $\dfrac1{10}-\dfrac1{12}=\dfrac1{60}$ ($60$ ngày).
\dapso{\pa{a} $10$ ngày; \pa{b} A: $30$, B: $20$, C: $60$ ngày.}}
\gap
\begin{probbox}{36}{\stars{2}}
Hai vòi cùng chảy vào một bể cạn. Cả hai vòi cùng chảy $2$ giờ thì được $\dfrac12$ bể. Vòi I chảy $3$ giờ và vòi II chảy $1$ giờ thì được $\dfrac7{12}$ bể. Hỏi mỗi vòi chảy một mình thì đầy bể trong bao lâu?
\end{probbox}
\probsol{
Trong $1$ giờ cả hai vòi chảy được $\dfrac12:2=\dfrac14$ bể.\par Trường hợp thứ hai gồm ``cả hai vòi chảy $1$ giờ'' và ``vòi I chảy thêm $2$ giờ''.\par
Vậy vòi I chảy được $\dfrac7{12}-\dfrac14=\dfrac13$ bể trong $2$ giờ, tức là\par $\dfrac16$ bể mỗi giờ.\par Vòi II: $\dfrac14-\dfrac16=\dfrac1{12}$ bể mỗi giờ.
\dapso{Vòi I: $6$ giờ; vòi II: $12$ giờ.}}
\gap
\begin{probbox}{37}{\stars{1}}
Một bể đang chứa nửa bể nước. Vòi I chảy đầy bể cạn trong $6$ giờ; vòi II tháo hết bể đầy trong $8$ giờ. Mở cả hai vòi thì sau bao lâu bể đầy?
\end{probbox}
\probsol{
Mỗi giờ: $\dfrac16-\dfrac18=\dfrac1{24}$ bể.\par Thời gian: $\dfrac12:\dfrac1{24}=12$ giờ.
\dapso{$12$ giờ.}}
\gap
\begin{probbox}{38}{\stars{1}}
A làm $3$ ngày và B làm $4$ ngày thì được $\dfrac5{12}$ công việc; A làm $6$ ngày và B làm $2$ ngày thì được $\dfrac7{12}$ công việc. Hỏi mỗi người làm một mình thì xong trong bao lâu?
\end{probbox}
\probsol{
Gấp đôi trường hợp thứ nhất: A làm $6$ ngày và B làm $8$ ngày thì được $\dfrac{10}{12}$ công việc.\par So với trường hợp thứ hai, $6$ ngày của B làm được $\dfrac{10}{12}-\dfrac7{12}=\dfrac14$ công việc, nên mỗi ngày B làm được $\dfrac1{24}$ công việc.\par $3$ ngày của A làm được $\dfrac5{12}-\dfrac4{24}=\dfrac14$, nên mỗi ngày A làm được $\dfrac1{12}$ công việc.\par A: $12$ ngày; B: $24$ ngày.
\dapso{A $12$, B $24$ ngày.}}
\gap
\begin{probbox}{39}{\stars{1}}
Hai người cùng làm thì xong một công việc trong $12$ giờ. Nếu người thứ nhất làm $4$ giờ và người thứ hai làm $6$ giờ thì được $\dfrac25$ công việc. Hỏi mỗi người làm một mình thì xong trong bao lâu?
\end{probbox}
\probsol{
Hai người cùng làm $4$ giờ thì được $\dfrac4{12}=\dfrac13$ công việc.\par $2$ giờ làm thêm của người thứ hai được $\dfrac25-\dfrac13=\dfrac1{15}$ công việc, tức là\par $\dfrac1{30}$ công việc mỗi giờ.\par Người thứ nhất: $\dfrac1{12}-\dfrac1{30}=\dfrac1{20}$ công việc mỗi giờ.\par Người thứ nhất: $20$ giờ; người thứ hai: $30$ giờ.
\dapso{$20$ giờ và $30$ giờ.}}
\gap
\begin{probbox}{40}{\stars{1}}
Một vòi chảy vào đầy bể trong $4$ giờ; một vòi tháo nước ra hết bể đầy trong $6$ giờ. Bể đang cạn, vòi chảy vào chảy $2$ giờ thì mở thêm vòi tháo. Hỏi sau bao lâu kể từ lúc bắt đầu thì bể đầy?
\end{probbox}
\probsol{
Sau $2$ giờ bể có $\dfrac24=\dfrac12$ bể nước.\par Khi mở cả hai vòi, mỗi giờ nước trong bể tăng thêm $\dfrac14-\dfrac16=\dfrac1{12}$ bể, cần thêm $\dfrac12:\dfrac1{12}=6$ giờ.\par Tất cả: $2+6=8$ giờ.
\dapso{$8$ giờ.}}
\gap
\begin{probbox}{41}{\stars{1}}
A và B cùng làm thì xong một công việc trong $6$ ngày; A làm một mình thì xong trong $10$ ngày. C làm một mình thì mất thời gian gấp đôi B. Hỏi cả ba người cùng làm thì xong trong bao lâu?
\end{probbox}
\probsol{
Mỗi ngày B làm được $\dfrac16-\dfrac1{10}=\dfrac1{15}$ công việc, nên B làm một mình cần $15$ ngày và C cần $30$ ngày.\par Cả ba người: $\dfrac1{10}+\dfrac1{15}+\dfrac1{30}=\dfrac6{30}=\dfrac15$ công việc mỗi ngày.\par Cùng làm: $5$ ngày.
\dapso{$5$ ngày.}}
\gap
\begin{probbox}{42}{\stars{1}}
Ba người cùng làm xong một công việc và được trả $3\,600\,000$ đồng. Nếu làm một mình thì A cần $6$ ngày, B cần $12$ ngày và C cần $4$ ngày. Chia tiền theo phần việc làm được thì mỗi người nhận được bao nhiêu?
\end{probbox}
\probsol{
Mỗi ngày A làm được $\dfrac2{12}$, B làm được $\dfrac1{12}$ và C làm được $\dfrac3{12}$ công việc; tỉ số là $2:1:3$.\par Một phần: $3\,600\,000:6=600\,000$ đồng.\par A: $1\,200\,000$; B: $600\,000$; C: $1\,800\,000$ đồng.
\dapso{A $1\,200\,000$; B $600\,000$; C $1\,800\,000$ đồng.}}
\gap
\begin{probbox}{43}{\stars{1}}
Cho vòi I chảy $2$ giờ rồi cho vòi II chảy $3$ giờ thì đầy bể. Cả hai vòi cùng chảy thì đầy bể trong $2$ giờ $24$ phút. Hỏi mỗi vòi chảy một mình thì đầy bể trong bao lâu?
\end{probbox}
\probsol{
$2$ giờ $24$ phút $=\dfrac{12}5$ giờ, nên mỗi giờ cả hai vòi chảy được $\dfrac5{12}$ bể.\par Cả hai vòi cùng chảy $2$ giờ được $\dfrac{10}{12}$ bể; một giờ chảy thêm của vòi II được phần còn lại, $1-\dfrac{10}{12}=\dfrac16$ bể.\par Vòi II: $6$ giờ; mỗi giờ vòi I chảy được $\dfrac5{12}-\dfrac16=\dfrac14$ bể: $4$ giờ.
\dapso{Vòi I $4$ giờ, vòi II $6$ giờ.}}
\gap
\begin{probbox}{44}{\stars{1}}
A làm một mình thì cần $10$ ngày, B làm một mình thì cần $15$ ngày. Hai người bắt đầu cùng lúc, nhưng A nghỉ $5$ ngày đầu, còn B làm liên tục. Hỏi công việc hoàn thành trong bao nhiêu ngày?
\end{probbox}
\probsol{
Trong $5$ ngày đầu B làm được $\dfrac5{15}=\dfrac13$ công việc; còn lại $\dfrac23$.\par Sau đó mỗi ngày cả hai làm được $\dfrac1{10}+\dfrac1{15}=\dfrac16$ công việc, cần $\dfrac23:\dfrac16=4$ ngày.\par Tất cả: $5+4=9$ ngày.
\dapso{$9$ ngày.}}
\gap
\begin{probbox}{45}{\stars{1}}
Một bể có hai vòi chảy vào và một vòi tháo ra. Vòi I chảy đầy bể trong $3$ giờ, vòi II trong $4$ giờ; vòi tháo làm cạn bể trong $6$ giờ. Bể đang cạn, mở cả ba vòi thì sau bao lâu bể đầy?
\end{probbox}
\probsol{
Mỗi giờ: $\dfrac13+\dfrac14-\dfrac16=\dfrac{4+3-2}{12}=\dfrac5{12}$ bể.\par Thời gian: $1:\dfrac5{12}=\dfrac{12}5$ giờ $=2$ giờ $24$ phút.
\dapso{$2$ giờ $24$ phút.}}
\gap
\begin{probbox}{46}{\stars{1}}
Lớp 5A làm vệ sinh sân trường xong trong $3$ giờ, lớp 5B xong trong $6$ giờ. Lớp 5A làm $1$ giờ, sau đó lớp 5B đến làm cùng. Hỏi hai lớp phải làm thêm bao lâu nữa?
\end{probbox}
\probsol{
Sau $1$ giờ lớp 5A làm được $\dfrac13$ công việc; còn lại $\dfrac23$.\par Mỗi giờ hai lớp cùng làm được: $\dfrac13+\dfrac16=\dfrac12$ công việc.\par Thời gian: $\dfrac23:\dfrac12=\dfrac43$ giờ $=1$ giờ $20$ phút.
\dapso{$1$ giờ $20$ phút.}}
\gap
\begin{probbox}{47}{\stars{1}}
Một máy bơm hút cạn nước một cái ao trong $10$ giờ. Sau $4$ giờ có thêm máy bơm thứ hai cùng bơm, và $2$ giờ sau thì ao cạn nước. Hỏi máy bơm thứ hai bơm một mình thì hút cạn ao trong bao lâu?
\end{probbox}
\probsol{
Trong $4$ giờ máy I hút được $\dfrac4{10}=\dfrac25$ lượng nước; còn lại $\dfrac35$.\par Hai máy hút $\dfrac35$ trong $2$ giờ: mỗi giờ $\dfrac3{10}$.\par Máy II: mỗi giờ $\dfrac3{10}-\dfrac1{10}=\dfrac15$, nên bơm một mình cần $5$ giờ.
\dapso{$5$ giờ.}}
\gap
\begin{probbox}{48}{\stars{2}}
Vòi I và vòi II cùng chảy thì đầy bể trong $2$ giờ; vòi II và vòi III cùng chảy thì đầy bể trong $3$ giờ; vòi I chảy một mình thì đầy bể trong $3$ giờ. Hỏi cả ba vòi cùng chảy thì đầy bể trong bao lâu?
\end{probbox}
\probsol{
Vòi II: mỗi giờ $\dfrac12-\dfrac13=\dfrac16$ bể.\par Vòi III: mỗi giờ $\dfrac13-\dfrac16=\dfrac16$ bể.\par Cả ba vòi: mỗi giờ $\dfrac13+\dfrac16+\dfrac16=\dfrac23$ bể.\par Thời gian: $1:\dfrac23=\dfrac32$ giờ $=1$ giờ $30$ phút.
\dapso{$1$ giờ $30$ phút.}}
\gap
\begin{probbox}{49}{\stars{2}}
Công nhân A làm một mình thì nhanh hơn công nhân B làm một mình $6$ giờ. Hai người cùng làm thì xong trong $4$ giờ. Hỏi mỗi người làm một mình thì xong trong bao lâu? (Gợi ý: thử chọn.)
\end{probbox}
\probsol{
Mỗi người làm một mình phải mất hơn $4$ giờ.\par Thử A $5$ giờ, B $11$ giờ: $\dfrac15+\dfrac1{11}=\dfrac{16}{55}\ne\dfrac14$.\par Thử A $6$ giờ, B $12$ giờ: $\dfrac16+\dfrac1{12}=\dfrac3{12}=\dfrac14$ (đúng).\par A: $6$ giờ; B: $12$ giờ.
\dapso{A $6$ giờ, B $12$ giờ.}}
\gap
\begin{probbox}{50}{\stars{2}}
Một người thợ cả mỗi ngày làm $8$ giờ thì làm xong một công việc trong $12$ ngày. Nếu có thêm một người thợ phụ có năng suất bằng một nửa người thợ cả cùng làm (mỗi ngày $8$ giờ) thì công việc đó hoàn thành trong bao nhiêu ngày?
\end{probbox}
\probsol{
Mỗi ngày người thợ cả làm được $\dfrac1{12}$ công việc, người thợ phụ làm được $\dfrac1{24}$ công việc.\par Cùng làm: $\dfrac1{12}+\dfrac1{24}=\dfrac18$ công việc mỗi ngày.\par Vậy cần $8$ ngày.
\dapso{$8$ ngày.}}
\clearpage
\titlehead
